
\begin{abstract}
光电容积脉搏波（photoplethysmography, PPG）可由腕表、指环等设备低成本获取，是无袖带血压估计最常用的信号源；PPG与血压存在统计关联，估计精度取决于表征形式、监督标签与验证协议是否匹配。本文以如何由 PPG 估计血压为主线：说明 PPG 的生理来源及其与血压的关联，再按手工特征与传统回归、端到端深度网络、预训练并微调三条路线，比较各类方法在表征、监督与部署代价上的取舍，并讨论个体化校准与多模态融合；随后整理常用数据集与临床评价标准，并讨论数据划分与校准隔离对结论可比性的影响。目前，血压估计的方法重心已由特征设计转向表征学习与个体校准，进一步提升的瓶颈则在于数据质量与评测协议；人群差异、校准负担与标签稀缺仍是走向实用的主要障碍。
\end{abstract}

\vspace{0.45\baselineskip}
